% GENERATED FILE. Edit canonical lessons, then run npm run generate:books.
% canonical-source-hash: fnv1a64:629e871ca01b8c1b
% canonical-lessons: IT-C332-tagliere, IT-C332-affettare, IT-C332-tritare, IT-C332-sbucciare, IT-C332-grattugiare

\chapter{Sul tagliere}
\label{ch:it-sul-tagliere}
\hlchaptermodality{332}

\begin{chapteropening}
\textbf{By the end of this chapter:} \emph{I can say how I prepared the vegetables: slicing, chopping, peeling and grating on the board.}

Complete the last lesson of chapter 332: I can say how I prepared the vegetables: slicing, chopping, peeling and grating on the board.
\end{chapteropening}

\section[il tagliere]{il tagliere — the chopping board}

\label{lesson:IT-C332-tagliere}



\begin{quote}
Before the new one: say the Italian for the teaspoon, then the Italian for the bowl.
\end{quote}

\subsection*{You'll want to know: il tagliere}
\textbf{il tagliere} — \textquotedblleft{}the chopping board\textquotedblright{}.

\begin{quote}
Taglia la cipolla sul tagliere.
\end{quote}

\begin{quote}
\emph{Cut the onion on the chopping board.}
\end{quote}

Built on \emph{tagliare}.

\begin{cousinweb}
\emph{tagliere} is built on \emph{tagliare} \textquotedblleft{}to cut\textquotedblright{}, from Late Latin \emph{tāliāre} \textquotedblleft{}to cut\textquotedblright{} (EtymDB).
\end{cousinweb}

\begin{grammarlens}[title={What you've built}]
One more word for telling, explaining and describing at B1.
\end{grammarlens}

\subsection*{Your turn}
\begin{itemize}
\raggedright
  \item \emph{Say it:} \emph{il tagliere}
  \item \emph{Say it:} \emph{il tagliere}, once more
  \item \emph{Say it:} a sentence of your own with \emph{il tagliere}
  \item \emph{Recall:} say the Italian for the teaspoon, then the Italian for the bowl, then say \emph{il tagliere} again
\end{itemize}

\begin{tcolorbox}[breakable,colback=teal!4,colframe=teal!35!black,title={Before you move on}]
What does il tagliere mean? (\textquotedblleft{}The chopping board\textquotedblright{}.) Say it once more.
\end{tcolorbox}
\section[affettare]{affettare — to slice}

\label{lesson:IT-C332-affettare}



\begin{quote}
Before the new one: say the Italian for the bowl, then the Italian for the chopping board.
\end{quote}

\subsection*{You'll want to know: affettare}
\textbf{affettare} — \textquotedblleft{}to slice\textquotedblright{}.

\begin{quote}
Ho affettato il pane e il prosciutto.
\end{quote}

\begin{quote}
\emph{I sliced the bread and the ham.}
\end{quote}

Built on \emph{fetta}, ‘slice’. \emph{Gli affettati} are cold cuts.

\begin{cousinweb}
\emph{affettare} is built on \emph{fetta} \textquotedblleft{}slice\textquotedblright{}, whose origin is not settled (EtymDB has no entry).
\end{cousinweb}

\begin{grammarlens}[title={What you've built}]
One more word for telling, explaining and describing at B1.
\end{grammarlens}

\subsection*{Your turn}
\begin{itemize}
\raggedright
  \item \emph{Say it:} \emph{affettare}
  \item \emph{Say it:} \emph{affettare}, once more
  \item \emph{Say it:} a sentence of your own with \emph{affettare}
  \item \emph{Recall:} say the Italian for the bowl, then the Italian for the chopping board, then say \emph{affettare} again
\end{itemize}

\begin{tcolorbox}[breakable,colback=teal!4,colframe=teal!35!black,title={Before you move on}]
What does affettare mean? (\textquotedblleft{}To slice\textquotedblright{}.) Say it once more.
\end{tcolorbox}
\section[tritare]{tritare — to chop finely, to mince}

\label{lesson:IT-C332-tritare}



\begin{quote}
Before the new one: say the Italian for the chopping board, then the Italian for to slice.
\end{quote}

\subsection*{You'll want to know: tritare}
\textbf{tritare} — \textquotedblleft{}to chop finely, to mince\textquotedblright{}.

\begin{quote}
Trita l'aglio e la cipolla.
\end{quote}

\begin{quote}
\emph{Chop the garlic and the onion finely.}
\end{quote}

\emph{La carne tritata} — mince.

\begin{cousinweb}
\emph{tritare} is built on \emph{trito} \textquotedblleft{}worn, ground\textquotedblright{}, from Latin \emph{trītus}, the past participle of \emph{terere} \textquotedblleft{}to rub\textquotedblright{} (Lewis \& Short).
\end{cousinweb}

\begin{grammarlens}[title={What you've built}]
One more word for telling, explaining and describing at B1.
\end{grammarlens}

\subsection*{Your turn}
\begin{itemize}
\raggedright
  \item \emph{Say it:} \emph{tritare}
  \item \emph{Say it:} \emph{tritare}, once more
  \item \emph{Say it:} a sentence of your own with \emph{tritare}
  \item \emph{Recall:} say the Italian for the chopping board, then the Italian for to slice, then say \emph{tritare} again
\end{itemize}

\begin{tcolorbox}[breakable,colback=teal!4,colframe=teal!35!black,title={Before you move on}]
What does tritare mean? (\textquotedblleft{}To chop finely, to mince\textquotedblright{}.) Say it once more.
\end{tcolorbox}
\section[sbucciare]{sbucciare — to peel}

\label{lesson:IT-C332-sbucciare}



\begin{quote}
Before the new one: say the Italian for to slice, then the Italian for to chop finely.
\end{quote}

\subsection*{You'll want to know: sbucciare}
\textbf{sbucciare} — \textquotedblleft{}to peel\textquotedblright{}.

\begin{quote}
Ho sbucciato le patate.
\end{quote}

\begin{quote}
\emph{I peeled the potatoes.}
\end{quote}

Built on \emph{buccia}, ‘peel, skin’.

\begin{cousinweb}
\emph{sbucciare} is \emph{s-} with \emph{buccia} \textquotedblleft{}peel\textquotedblright{}, whose origin is not settled (EtymDB has no entry).
\end{cousinweb}

\begin{grammarlens}[title={What you've built}]
One more word for telling, explaining and describing at B1.
\end{grammarlens}

\subsection*{Your turn}
\begin{itemize}
\raggedright
  \item \emph{Say it:} \emph{sbucciare}
  \item \emph{Say it:} \emph{sbucciare}, once more
  \item \emph{Say it:} a sentence of your own with \emph{sbucciare}
  \item \emph{Recall:} say the Italian for to slice, then the Italian for to chop finely, then say \emph{sbucciare} again
\end{itemize}

\begin{tcolorbox}[breakable,colback=teal!4,colframe=teal!35!black,title={Before you move on}]
What does sbucciare mean? (\textquotedblleft{}To peel\textquotedblright{}.) Say it once more.
\end{tcolorbox}
\section[grattugiare]{grattugiare — to grate}

\label{lesson:IT-C332-grattugiare}



\begin{quote}
Before the new one: say the Italian for to chop finely, then the Italian for to peel.
\end{quote}

\subsection*{You'll want to know: grattugiare}
\textbf{grattugiare} — \textquotedblleft{}to grate\textquotedblright{}.

\begin{quote}
Grattugia un po' di formaggio sulla pasta.
\end{quote}

\begin{quote}
\emph{Grate a little cheese over the pasta.}
\end{quote}

Present \emph{grattugio}, \emph{grattugi}.

\begin{cousinweb}
\emph{grattugiare} is built on \emph{grattugia} \textquotedblleft{}grater\textquotedblright{}, from \emph{grattare} \textquotedblleft{}to scratch\textquotedblright{}, which is usually traced to a Germanic word (EtymDB has no entry).
\end{cousinweb}

\begin{grammarlens}[title={What you've built}]
One more word for telling, explaining and describing at B1.
\end{grammarlens}

\subsection*{Your turn}
\begin{itemize}
\raggedright
  \item \emph{Say it:} \emph{grattugiare}
  \item \emph{Say it:} \emph{grattugiare}, once more
  \item \emph{Say it:} a sentence of your own with \emph{grattugiare}
  \item \emph{Recall:} say the Italian for to chop finely, then the Italian for to peel, then say \emph{grattugiare} again
\end{itemize}

\begin{tcolorbox}[breakable,colback=teal!4,colframe=teal!35!black,title={Before you move on}]
What does grattugiare mean? (\textquotedblleft{}To grate\textquotedblright{}.) Say it once more.
\end{tcolorbox}
